% Generated draft: requires mathematical and source review.
\documentclass[11pt]{article}
\usepackage[margin=25mm]{geometry}
\usepackage{fontspec}
\setmainfont{Latin Modern Roman}
\usepackage{amsmath,amssymb,mathtools,unicode-math}
\setmathfont{Latin Modern Math}
\usepackage{enumitem,needspace,xurl,hyperref}
\setlength{\parindent}{0pt}
\setlength{\parskip}{5pt}
\setlength{\emergencystretch}{3em}
\newcommand{\N}{\mathbb{N}}
\newcommand{\Z}{\mathbb{Z}}
\newcommand{\Q}{\mathbb{Q}}
\newcommand{\R}{\mathbb{R}}
\newcommand{\C}{\mathbb{C}}
\newcommand{\F}{\mathbb{F}}
\newcommand{\A}{\mathbb{A}}
\newcommand{\PP}{\mathbb{P}}
\newcommand{\E}{\mathbb{E}}
\newcommand{\eps}{\varepsilon}
\newcommand{\dd}{\,\mathrm{d}}
\begin{document}
\section*{UnsolvedMath ID: 10600029}
\textbf{Source problem code:} catalogue entry
\subsection{Definitions and mathematical statement}
In a Khovanov cube of resolutions, each classical crossing is smoothed in two ways and edges correspond to changing one smoothing. Ordinary merge and split saddles act through multiplication and comultiplication in the relevant Frobenius algebra. In a virtual diagram, an edge may instead change one state circle into one state circle. The source identifies this one-to-one resmoothing as an integral-coefficient issue, while reporting the mod-2 construction straightforward.

The task is to explain the integral virtual-knot/link construction in terms of Khovanov's original framework, using the source's claimed connection with Khovanov--Rozansky categorification of infinitely many HOMFLYPT specializations, so as to clarify Manturov's integral construction. An explicit comparison must identify the relevant specialization, state objects, local maps, signs and grading conventions, including the one-to-one resmoothing case, and justify the claimed integral homology and virtual Reidemeister invariance.

The source already reports an integral construction by Manturov and says the Khovanov--Rozansky approach solves the integral problem in principle. Thus it asks for the technical bridge/clarification, not an unqualified new existence theorem. The excerpt does not identify the exact complexes or specify whether comparison means chain isomorphism, chain homotopy equivalence or homology isomorphism. No equality of all HOMFLYPT specializations is asserted, and passing only to mod-2 or rational coefficients would omit the integral issue and possible torsion.

The summary calls the problem solved and the dictionary largely accomplished, but supplies no actual comparison theorem or references. Its conclusion is retained as a status report requiring verification, not adopted as a checked resolution of the requested technical bridge.
\subsection{Short English statement}
Make the integral virtual Khovanov construction explicit in the original cube-of-resolutions language, especially one-to-one resmoothings, and explain its relation to the stated Khovanov--Rozansky approach and Manturov's construction.
\subsection{Sources}
\begin{itemize}
\item Local UnsolvedMath registry, permanent ID 10600029, source code catalogue entry; no proposer or proposal year supplied; updated 2026-08-06; classification and status fields say partial progress, while the summary calls the technical problem solved.
\item \emph{Unsolved Problems in Virtual Knot Theory and Combinatorial Knot Theory}, revised survey problem-list item 29: \url{https://arxiv.org/abs/1409.2823}. Registry access date 2026-07-29; source-TeX extraction. Provenance reports 60 top-level items, with compound items retained. Present primary-literature status and redistribution terms require review.
\item Manturov and Khovanov--Rozansky are named in the source, but their citation slots/bibliographic details are absent. The summary additionally mentions KLS-type constructions without defining the abbreviation or providing an exact reference. No citation is reconstructed here.
\end{itemize}
\end{document}
